%%% FIELD hoeffding-statement
Let \(N\ge 1\), and let \(Y_1,\ldots,Y_N\) be independent real-valued random variables such that \(a_i\le Y_i\le b_i\) almost surely.  Put \(\overline Y=N^{-1}\sum_{i=1}^N Y_i\) and \(\mu=\mathbb E[\overline Y]\).  Then, for every \(u>0\),
\[
 \Pr\{\overline Y-\mu\ge u\}
 \le
 \exp\!\left\{-\frac{2N^2u^2}{\sum_{i=1}^N(b_i-a_i)^2}\right\}.
\]
%%% FIELD target-statement
Suppose each transition row has \(n_{x,a}\) independent multinomial calibration draws, the nominal zero pattern is known, and calibration is independent of the factual record. Then
\[
\Pr_{\mathrm{cal}}\left\{I_T(P,\rho,b,o,e)\subseteq\widehat I_{T,\alpha}(P,\rho,b,o,e)\text{ for every }T\ge1\text{ and every }e\in\mathcal O^{T+1}\text{ with }\Pr_{P,\rho,b,o}(E_T=e)>0\right\}\ge1-\alpha.
\]
Both coarsened-record endpoints are therefore covered simultaneously for every fixed factual-law-positive record, even when response masses are zero and the endpoint-active face changes.
%%% FIELD target-proof
For each \((x,a)\in\mathcal S\times\mathcal A\), write \(Y^{x,a}_1,\ldots,Y^{x,a}_{n_{x,a}}\) for the multinomial calibration draws from the row \(P(\cdot\mid x,a)\).  For \(y\in\mathcal S\), define
\[
 V^{x,a,y}_i=\mathbf 1\{Y^{x,a}_i=y\},
 \qquad
 \widehat P(y\mid x,a)=\frac{1}{n_{x,a}}\sum_{i=1}^{n_{x,a}}V^{x,a,y}_i.
\]
The variables \(V^{x,a,y}_i\), with \(i\) varying and \((x,a,y)\) fixed, are independent, take values in \([0,1]\), and have common expectation \(P(y\mid x,a)\), by the stated rowwise multinomial sampling scheme.  Since \(d\ge2\), \(k\ge1\), \(0<\alpha<1\), and \(n_{x,a}\ge1\), the radius \(\delta_{x,a}\) is strictly positive.  Applying Lemma \(\mathrm{lem:hoeffding\mbox{-}one\mbox{-}sided}\) first to the variables \(V^{x,a,y}_i\) and then to \(-V^{x,a,y}_i\), and adding the two one-sided bounds, gives
\[
 \Pr_{\mathrm{cal}}\!\left\{
 \left|\widehat P(y\mid x,a)-P(y\mid x,a)\right|>\delta_{x,a}
 \right\}
 \le 2\exp\{-2n_{x,a}\delta_{x,a}^{,2}\}
 =\frac{\alpha}{d^2k}.
\]
No independence between distinct cells is required for the next step.  Define the single calibration event
\[
 \mathcal H=
 \bigcap_{x\in\mathcal S}\bigcap_{a\in\mathcal A}\bigcap_{y\in\mathcal S}
 \left\{
 \left|\widehat P(y\mid x,a)-P(y\mid x,a)\right|\le\delta_{x,a}
 \right\}.
\]
There are exactly \(d^2k\) indexed cells.  Therefore the union bound and the preceding display imply
\[
 \Pr_{\mathrm{cal}}(\mathcal H^c)
 \le
 \sum_{x\in\mathcal S}\sum_{a\in\mathcal A}\sum_{y\in\mathcal S}
 \frac{\alpha}{d^2k}
 =\alpha,
 \qquad
 \Pr_{\mathrm{cal}}(\mathcal H)\ge1-\alpha.
\]

We next show deterministically that \(\mathcal H\) implies the asserted simultaneous inclusion.  On \(\mathcal H\), the true kernel \(P\) obeys every cellwise inequality in Definition \(\mathrm{def:confidence\mbox{-}region}\).  It has the stipulated zero pattern because that pattern is the nominal one.  On each nonzero cell it obeys \(P(y\mid x,a)\ge\epsilon\) by Assumption \(\mathrm{ass:positive\mbox{-}cells}\).  Finally, the pair \((g,h)\) witnesses the last two constraints in that definition: Assumption \(\mathrm{ass:poisson\mbox{-}equation}\) gives
\[
 r-g\mathbf 1=h-K_{\pi}h,
\]
and Assumption \(\mathrm{ass:poisson\mbox{-}span}\) gives \(\operatorname{span}(h)\le B\).  Hence
\[
 P\in\mathcal C_{\alpha}
 \quad\text{on }\mathcal H,
\]
so in particular \(\mathcal C_{\alpha}\ne\varnothing\) there.

Fix any integer \(T\ge1\) and any \(e\in\mathcal O^{T+1}\) satisfying \(\Pr_{P,\rho,b,o}(E_T=e)>0\).  We record why all intervals in the defining union are well typed.  By Definition \(\mathrm{def:factual\mbox{-}record\mbox{-}law}\), positivity of this probability implies, because the displayed sum is finite and has nonnegative summands, that there is a state path \(s_{0:T}\) such that \(\rho(s_0)>0\), \(o(s_t)=e_t\) for every \(t\), and
\[
 P(s_t\mid s_{t-1},b(s_{t-1}))>0
 \quad (t=1,\ldots,T).
\]
Every \(P'\in\mathcal C_{\alpha}\) has the same zero pattern as \(P\); consequently each corresponding transition probability under \(P'\) is positive.  The same path then contributes a strictly positive summand to \(\Pr_{P',\rho,b,o}(E_T=e)\).  Thus \(I_T(P',\rho,b,o,e)\) is defined for every retained \(P'\), as required by Definition \(\mathrm{def:sharp\mbox{-}interval}\).

On \(\mathcal H\), the membership \(P\in\mathcal C_{\alpha}\) and Definition \(\mathrm{def:confidence\mbox{-}fiber}\) now give, simultaneously for this and every other admissible \((T,e)\),
\[
 \begin{aligned}
 I_T(P,\rho,b,o,e)
 &\subseteq
 \bigcup_{P'\in\mathcal C_{\alpha}}I_T(P',\rho,b,o,e)\\
 &\subseteq
 \operatorname{hull}\!\left(
   \bigcup_{P'\in\mathcal C_{\alpha}}I_T(P',\rho,b,o,e)
 \right)
 =\widehat I_{T,\alpha}(P,\rho,b,o,e).
 \end{aligned}
\]
The event \(\mathcal H\) does not depend on \(T\), on \(e\), or on an optimizer over \(\Theta(P)\).  It is therefore contained in the event inside the probability in the theorem.  The bound \(\Pr_{\mathrm{cal}}(\mathcal H)\ge1-\alpha\) proves the displayed simultaneous claim.

For completeness, the independence assumption supplies the advertised post-record interpretation rather than being needed for the preceding set inclusion.  The event \(\mathcal H\) is measurable with respect to the calibration sample.  By Assumption \(\mathrm{ass:calibration\mbox{-}independence}\), for every fixed \(T\) and every factual-law-positive \(e\),
\[
 \Pr\{\mathcal H\mid E_T=e\}
 =\Pr_{\mathrm{cal}}(\mathcal H)
 \ge1-\alpha.
\]
Since \(\mathcal H\) implies inclusion of the entire sharp interval, this is simultaneous conditional coverage of its lower and upper endpoints.  The argument never divides by a response-program mass, assumes that such a mass is positive, differentiates an endpoint optimizer, or requires the optimizer to remain on one face of \(\Theta(P)\).
